\chapter{Rappels et notation}

Soit $V$ un espace vectoriel sur un corps de dimension finie $K$.

\begin{notation}[Dimension d'un espace vectoriel]
$$
\dim(V)=n<\infty
$$
\end{notation}

\begin{definition}[Base de $V$]
On la note $\mathscr{B}=\{ v_{1},\dots,v_{n} \}$
\begin{itemize}
\item Une famille libre
\item Tout $v\in V$ est une combinaison lineaire qu'on exprime:
$$
v=v_{1}\alpha_{1}+\dots+\alpha _{n}v_{n}
$$
$$
[v]_{\mathscr{B}}=\begin{pmatrix}
\alpha_{1} \\
\vdots \\
\alpha _{n} 
\end{pmatrix}\in K^{n}
$$
\end{itemize}
\end{definition}
\begin{notation}[Isomorphisme]
\begin{align*}
\psi: V & \to K^{n} \\
v & \mapsto[v]_{\mathscr{B}}
\end{align*}
\end{notation}

\begin{definition}[Changement de base]
Soit $\mathscr{B}'=\{ v_{1},\dots,v_{n} \}$ une (autre) base de $V$. Il existe une unique matrice $P_{\mathscr{B}\mathscr{B}'}\in K^{n\times n}$ inversible telle que  pour tout $v\in V$:
$$
[v]_{\mathscr{B'}}=P_{\mathscr{B}\mathscr{B}'}[v]_{\mathscr{B}}
$$
\end{definition}
\begin{example}
Soit $V$ de $\dim(V)=3$ avec une base $\mathscr{B}=\{ v_{1},v_{2},v_{3} \}$ et une autre base $\mathscr{B}'=\{ w_{1},w_{2},w_{3} \}$. On a les relations
\begin{align*}
w_{1} & =2v_{1}+v_{2}-v_{3} \\
w_{2} & =v_{1}+3v_{2}+v_{3} \\
w_{3} & =v_{1}+v_{2}-v_{3}
\end{align*}
donc la matrice de changement:
$$
P_{\mathscr{B}'\mathscr{B}}=\begin{pmatrix}
2 & 1 & 1 \\
1 & 3 & 1 \\
-1 & 1 & -1 
\end{pmatrix}
$$
\end{example}
\begin{notation}
$$
\boldsymbol{\ell}_{1}=\begin{pmatrix}
1 \\ 0\\0
\end{pmatrix}
$$
\end{notation}
\begin{remark}
$P_{\mathscr{B}\mathscr{B}'}=P_{\mathscr{B}'\mathscr{B}}^{-1}$
\end{remark}
\section{Applications lineaires}
Soient $V,W$ deux espaces vectoriels sur $K$. $f:V\to W$ est une application lineaire si pour tout $u,v\in V$ et $\alpha \in K$ :
\begin{itemize}
\item $f(u+v)=f(u)+f(v)$
\item $f(\alpha u)=\alpha f(u)$
\end{itemize}
Soient $\dim(V)=n$, $\dim(W)=m$ et $\mathscr{B}_{V}=\{ v_{1}$,$\dots,v_{n} \}$, $\mathscr{B}_{W}=\{ w_{1},\dots,w_{m} \}$ les bases de $V$ et $W$ respectivement. $f$ est decrite par les images $$f(v_{j})=a_{1j}w_{1}+\dots+a_{mj}w_{m}$$ ou les coefficients $a_{ij}$ donne les coordones de $f(v_{i})$ dans $\mathscr{B}_{W}$.
\begin{example}
$$
f(2v_{1}+3v_{2})=2f(v_{1})+3f(v_{2})
$$
\end{example}

On peut construire ainsi la matrice associe a $f$
$$
\begin{pmatrix}
a_{11} & \dots & a_{1n} \\
\vdots &  & \vdots \\
a_{m1} & \dots & a_{mn}
\end{pmatrix}[v]_{\mathscr{B}_{V}}=[f(v)]_{\mathscr{B}_{W}}
$$
\begin{example}
Soien $V,W$ des espaces vectoriel sur $\mathbf{Z}_{3}$ des bases $\mathscr{B}_{V}=\{ v_{1},v_{2},\dots,v_{5} \}$ et $\mathscr{B}_{W}=\{ w_{1},w_{2},w_{3},w_{4}\}$. On a les images:
\begin{align*}
f(v_{1}) & =1w_{1}+2w_{2}+2w_{3}+2w_{4} \\
f(v_{2}) & =2w_{1}+w_{2}+w_{3}+w_{4} \\
f(v_{3}) & =2w_{2}+2w_{4}\\
f(v_{4}) & =w_{1}+2w_{2}+2w_{3}+2w_{4} \\
f(v_{5}) & =2w_{1}+2w_{2}+w_{3}+2w_{4}
\end{align*}
donc on obtient la matrice
$$
A=\begin{pmatrix}
1 & 2 & 0 & 1 & 2 \\
2 & 1 & 2 & 2 & 2 \\
2 & 1 & 0 & 2 & 1 \\
2 & 1 & 2 & 2 & 2
\end{pmatrix}\in \mathbf{Z}_{3}^{4\times 5}
$$
on a que $\forall v\in V$ : $[f(v)]_{\mathscr{B}_{W}}=A[v]_{\mathscr{B}_{V}}$.

Pour obtenir de bases de $\ker(f)$ et $\mathrm{Im}(f)$ on cherche une matrice $Q$ telle que $QA$ echelonne $A$ :
$$
QA=\begin{pmatrix}
1 & 2 & 0 & 1 & 2 \\
0 & 0 & 2 & 0 & 1 \\
0 & 0 & 0 & 0 & 0 \\
0 & 0 & 0 & 0 & 0
\end{pmatrix}
$$
donc on voit que $\mathrm{Im}(f)=\{ f(v_{1}),f(v_{3}) \}$ (les colonnes de pivots). Pour trouver la base de $\ker(f)$ on cherche les combinaisons lineaires des colonnes egales a 0 et on trouve $(e_{1}+e_{3}+e_{5},e_{4}+2e_{1},e_{2}+e_{1})$ ou $e_{i}$ est une colonne qu'on associe a $v_{i}$ pour avoir $\ker(f)=\{ v_{1}+v_{3}+v_{5},v_{4}+2v_{1},v_{2}+v_{1} \}$. 
\end{example}
\begin{definition}[Image]
$\mathrm{Im}(f)=\{ f(v):v\in V \}\subseteq W$
\end{definition}
\begin{definition}[Noyau]
$\ker(f)=\{ v\in V:f(v)=0\}\subseteq V$
\end{definition}
\begin{theorem}[Noyau-Image]
$$
\dim(\mathrm{Im}(f))+\dim(\ker(f))=\dim(V)
$$
\end{theorem}

\subsection{Comment trouver des bases de l'image et noyau}

\begin{definition}[Algorithme de Gauss]
Soit $A\in K^{m\times n}$, on cherche une matrice $Q\in K^{m\times m}$ inversible construit comme un produit des matrices des transformations elementaires telle que $QA$ soit une matrice échelonne qui aura $k$ pivots qu'on note $p_{i},1\leq i\leq k$. 
\end{definition}

\section{Endomorphismes et matrices semblables}

\begin{definition}[Endomorphisme]
Un endomorphisme est une application lineaire $f:V\to V$ de $V$ un espace vectoriel de $\dim(V)=n$ et base $\mathscr{B}=\{ v_{1},\dots,v_{n} \}$. 
\end{definition}
\begin{notation}
Pour tout $v\in V$ et $f$ une application lineaire on a:
$$
[f(v)]_{\mathscr{B}}=A_{\mathscr{B}}^{f}[v]_{\mathscr{B}}
$$
\end{notation}

Soit $\mathscr{B}'$ une autre de base de $V$ et $P_{\mathscr{B}'\mathscr{B}}$ une matrice de changement on a
$$
A^{f}_{\mathscr{B}'}=P_{\mathscr{B}'\mathscr{B}}^{-1}A_{\mathscr{B}}^{f}P_{\mathscr{B}'\mathscr{B}}
$$

\begin{definition}[Matrices semblables]
Soient $A,B\in K^{n\times n}$ on dit que ces matrices sont semblables $A\sim B$ s'il existe $P\in K^{n\times n}$ inversible telle que 
$$
A=P^{-1}BP
$$
\end{definition}

Une matrice $A\in K^{n\times n}$ est diagonalisable si
$$
A\sim \begin{pmatrix}\lambda_{1} &  &   \\  & \ddots &  \\0  &  & \lambda _{n}\end{pmatrix}\in K^{n\times n}
$$
 
\begin{question}
    Quand $A$ est diagonalisable?
\end{question}
\begin{answer}
    Si et seulement si le polynome minimale $p(x)\in K[X]$ de $A$ se factorise comme $p(x)=(x-\lambda_{1})\dots(x-\lambda _{k})$ avec les $\lambda _{i}$ distincts.  
\end{answer}

\begin{exercise}[Serie 1 exo plus]
Soit $R$ un anneau et $\alpha \in Z(R)$ un element du centre de $R$. Montrer que l'application

\begin{align}
\phi : R[x] & \to R \\
f(x) & \mapsto f(\alpha)
\end{align}

est un morphisme d'anneuax surjectif.

Soit $R=\mathbf{Z}^{2\times 2}$, trouver $\alpha \in R$ (qui n'appartient donc pas a $Z(R)$) et deux polynomes $g,h\in R[x]$ tel que
$$
g(\alpha)h(\alpha)\neq (gh)(\alpha)
$$
\end{exercise}
\begin{proof}[Solution]
Pour montrer que $\phi$ est un homomorphisme d'anneaux on verifie ces trois conditions $\forall f,g\in R[x]$
\begin{enumerate}
\item $\phi(1)=1$
\item $\phi(g+f)=\phi(g)+\phi(f)$
\item $\phi(fg)=\phi(f)\phi(g)$
\end{enumerate}
on verifie:
\begin{enumerate}
\item $\phi(1)=1$ car le polynome constant est envoye sur lui meme par $\phi$
\item \label{exo1} On prend $f,g\in R[x]$ et on ecrit $f(x)=\sum_{i=0}^{\deg f}a_{i}x^{i}$ et $g(x)=\sum_{j=0}^{\deg g}b_{j}x^{j}$ ou $a_{i},b_{j}\in R$. 
\begin{align}
\phi(f(x)+g(x)) & =\phi\left( \sum_{i=0}^{\deg f} a_{i}x^{i}+\sum_{j=0}^{\deg g} b_{j}x^{j}  \right) \\
 & =\phi\left( \sum_{k=0}^{\max\{ \deg f,\deg g \}} (a_{k}+b_{k})x^{k} \right) \\
 & =\sum_{k=0}^{\max\{ \deg f,\deg g \}}(a_{k}+b_{k})\alpha^{k} \\
 & =\sum_{i=0}^{\deg f} a_{i}\alpha^{i}+\sum_{j=0}^{\deg g} b_{j}\alpha^{j}  \\
 & =\phi(f(x))+\phi(g(x)) \\
\end{align}
\item On reprend $f,g$ comme avant et on calcule 
\begin{align}
\phi(f(x)g(x)) & =\phi\left( \sum_{i=0}^{\deg f} a_{i}x^{i} \sum_{j=0}^{\deg g} b_{j}x^{j} \right) \\
 & =\phi\left( \sum_{i=0}^{\deg f} \sum_{j=0}^{\deg g} a_{i}b_{j}x^{i+j} \right) \\
 & =\sum_{i=0}^{\deg f} \sum_{j=0}^{\deg g} a_{i}b_{j}\alpha^{i+j} \label{exo2}\\
 & =\left( \sum_{i=0}^{\deg f } a_{i}\alpha^{i} \right)\left( \sum_{j=0}^{\deg g } b_{j}\alpha^{j} \right) \\
 & =\phi(f(x))\phi(g(x))
\end{align}
\end{enumerate}
ou \ref{exo2} est obtenu car $\alpha \in Z(R)$. Donc $\phi$ est un morphisme d'anneaux qui est surjectif a cause de \ref{exo1}.

Soit $R=\mathbf{Z}^{2\times2}$, on prend $\alpha=\begin{pmatrix}1 & 1 \\ 0 & 1\end{pmatrix}$, $g(x)=\begin{pmatrix}1 & 1  \\ 0 & 1 \end{pmatrix}x$ et $h(x)=\begin{pmatrix} 1 & 0 \\ 1 & 1\end{pmatrix}x$ on a
$$
\begin{pmatrix}
1 & 1 \\
0 & 1
\end{pmatrix}\begin{pmatrix}
1 & 0 \\
1 & 1
\end{pmatrix}=\begin{pmatrix}
2 & 1 \\
1 & 1
\end{pmatrix}
$$
et 
$$
\begin{pmatrix}
1 & 0 \\
1 & 1
\end{pmatrix}\begin{pmatrix}
1 & 1 \\
0 & 1
\end{pmatrix}=\begin{pmatrix}
1 & 1 \\
1 & 2
\end{pmatrix}
$$
Donc ces matrices ne commutent pas. On obtient
$$
g(\alpha)=\begin{pmatrix}
1 & 2 \\
0 & 1
\end{pmatrix}\quad f(\alpha)=\begin{pmatrix}
1 & 1 \\
1 & 2
\end{pmatrix}
$$
donc
$$
g(\alpha)f(\alpha)=\begin{pmatrix}
3 & 5 \\
1 & 2
\end{pmatrix}
$$
mais
$$
g(x)h(x)=\begin{pmatrix}
2 & 1 \\
1 & 1 
\end{pmatrix}x^{2}
$$
que si on evalue en $\alpha$
$$
(gh)(\alpha)=\begin{pmatrix}
2 & 5 \\
1 & 3
\end{pmatrix}
$$
qui est different de $g(\alpha)f(\alpha)$
\end{proof}
